% !TeX program = lualatex
% Compile twice: lualatex 182.tex
% All references and prose are embedded; no BibTeX database or external inputs.
\documentclass[11pt,a4paper]{article}
\usepackage[margin=24mm,headheight=14pt]{geometry}
\usepackage{fontspec}
\setmainfont{Latin Modern Roman}
\setsansfont{Latin Modern Sans}
\setmonofont{Latin Modern Mono}
\usepackage{amsmath,amssymb,mathtools}
\usepackage{unicode-math}
\setmathfont{Latin Modern Math}
\usepackage{microtype}
\usepackage{enumitem,xurl,needspace}
\usepackage[dvipsnames]{xcolor}
\usepackage{hyperref}
\hypersetup{unicode=true,colorlinks=true,linkcolor=MidnightBlue,urlcolor=MidnightBlue,
 pdftitle={500 Mathematical Problem Statements},pdfauthor={Source-annotated compilation},bookmarksnumbered=true}
\usepackage{bookmark}
\usepackage{tocloft}
\setlength{\cftsecnumwidth}{3em}
\cftsetpnumwidth{2.5em}
\cftsetrmarg{4em}
\usepackage{titlesec}
\titleformat{\section}{\Large\bfseries\raggedright}{\thesection}{0.7em}{}
\usepackage{fancyhdr}
\pagestyle{fancy}\fancyhf{}
\fancyhead[L]{\small\sffamily Mathematical problem statements}
\fancyhead[R]{\small Source edition 22}
\fancyfoot[C]{\thepage}
\setlength{\parindent}{0pt}
\setlength{\parskip}{5pt plus 1pt}
\setlength{\emergencystretch}{3em}
\setcounter{tocdepth}{1}
\setcounter{secnumdepth}{1}
\setlist[itemize]{leftmargin=3em,itemsep=5pt,topsep=4pt,parsep=0pt}
\allowdisplaybreaks
\newcommand{\N}{\mathbb{N}}
\newcommand{\Z}{\mathbb{Z}}
\newcommand{\Q}{\mathbb{Q}}
\newcommand{\R}{\mathbb{R}}
\newcommand{\C}{\mathbb{C}}
\newcommand{\F}{\mathbb{F}}
\newcommand{\A}{\mathbb{A}}
\newcommand{\PP}{\mathbb P}
\newcommand{\E}{\mathbb{E}}
\newcommand{\eps}{\varepsilon}
\newcommand{\dd}{\,\mathrm d}
\newcommand{\sref}[2]{\hyperlink{source:#1:#2}{[S#2]}}
\newcommand{\eref}[2]{\hyperlink{extra:#1:#2}{[E#2]}}
% Turn the next switch off for a shorter reading copy. The quotations preserve
% every exactTarget field, including qualifications omitted from a short summary.
\newif\ifshowcatalogue
\showcataloguetrue
\newcommand{\cataloguescope}[1]{\ifshowcatalogue
 \paragraph{Exact scope in the supplied catalogue.}
 \begin{quote}\small #1\end{quote}\fi}

% Standalone notebook wrapper; original problem section retained verbatim.
\numberwithin{equation}{section}
\hypersetup{pdftitle={Research attempt -- problem 182},
 pdfauthor={Research notebook; source compilation retained}}
\fancyhead[L]{\small\sffamily Research notebook}
\fancyhead[R]{\small\sffamily Problem 182 | 27 September 2026}
\begin{document}
\setcounter{section}{181}
% ================================================================
% problemId: problem.voevodsky-smash-nilpotence-conjecture
% source releaseRank: 182; source releaseStatus: open_with_solved_subcases
% exactTarget SHA-256: c218d23023386e9060e34781b4560c08d126d2f6983d479159a7aeed36e8d592
\clearpage
\section{\texorpdfstring{Voevodsky's Smash-Nilpotence Conjecture}{Voevodsky's Smash-Nilpotence Conjecture}}\label{182}
\subsection{Definitions and mathematical statement}
For a smooth projective $d$-dimensional variety over $k$, a rational cycle class $\alpha\in\mathrm{CH}^r(X)_{\Q}$ is numerically trivial when
$\deg(\alpha\cdot\beta)=0$ for every complementary class $\beta\in\mathrm{CH}^{d-r}(X)_{\Q}$. It is smash-nilpotent when
\[
 \alpha^{\boxtimes N}=0\quad\text{in}\quad\mathrm{CH}^{Nr}(X^N)_{\Q}
 \quad\text{for some }N\ge1.
\]
Here $\boxtimes$ is the external product over $k$, and the exponent may depend on $\alpha$. The conjecture says numerical triviality implies smash-nilpotence for every field and every such cycle. Powers under intersection on $X$, or under composition of correspondences, are different operations. Rational coefficients and the absence of a restriction on the characteristic of $k$ are essential scope choices.
\par\smallskip\textit{Formulation and concept references:} \sref{182}{1}.
\cataloguescope{For every field k, every smooth projective k-variety X of dimension d, every integer 0<=r<=d, and every class alpha in CH\textasciicircum{}r(X)\_Q (codimension-r cycles modulo rational equivalence with rational coefficients), assume deg(alpha.beta)=0 for every beta in CH\textasciicircum{}(d-r)(X)\_Q. Voevodsky's smash-nilpotence assertion is that there exists an integer N>=1, depending on alpha, for which the N-fold EXTERNAL product alpha times ... times alpha is zero in CH\textasciicircum{}(Nr)(X\textasciicircum{}N)\_Q, where X\textasciicircum{}N is the fibre product over k. No algebraic-closure, characteristic-zero or geometric-connectedness hypothesis is imposed; disconnected smooth projective schemes are interpreted on their pure-dimensional components. The reverse implication, smash-nilpotent implies numerically trivial, is known, so this is equality of numerical and smash equivalence for rational cycles. Neither intersection powers on X, composition powers of correspondences, integral coefficients, nor a uniform or computable exponent replaces this assertion.}
\subsection{Short English statement}
Does every rational algebraic cycle invisible to all intersection-number tests become zero after taking sufficiently many external copies of itself?
% BEGIN RESEARCH ADDITION ATTEMPT-182
\subsection{Research attempt}

\paragraph{Current frontier.}
The foundational theorem is stronger than a statement about intersection
powers but weaker than the conjecture: Voevodsky proves smash-nilpotence
for cycles \emph{algebraically} equivalent to zero. His Proposition 3.1
handles degree-zero cycles on smooth projective curves, and Corollary 3.2
passes to general algebraically trivial cycles \sref{182}{1}.
Numerical triviality is not the hypothesis of that theorem. A 2026
manuscript of Corradini proves a tropical algebraic-triviality analogue
and formulates a tropical smash-nilpotence conjecture \eref{182}{1}.
It concerns a different category of cycles and supplies no theorem for
all numerically trivial Chow classes over arbitrary fields.

\paragraph{Attack route.}
One route would compare numerical and algebraic equivalence after taking
an external power. Another would seek tensor control from finite-dimensional
motives; composition nilpotence alone is tested below and found inadequate.
A third issue, often hidden in both routes, is descent from an algebraic
closure. We combine explicit exponent bookkeeping with descent, then
isolate the genuinely missing geometric input. All tensor products in the
argument are external products over the specified field.

\paragraph{Attempt: descent preserves the exact target.}
\textbf{Descent lemma.} If the conjecture is established over every
algebraically closed field, it follows over every field, with rational
coefficients as in the question.

\textit{Proof.} Let $\alpha$ be numerically trivial on $X/k$. For a finite
extension $L/k$, write $p:X_L\to X$. Every complementary cycle $\beta$ on
$X_L$ satisfies
\[
 [L:k]\deg_L(\alpha_L\cdot\beta)
   =\deg_k p_*(\alpha_L\cdot\beta)
   =\deg_k(\alpha\cdot p_*\beta)=0.
\]
This is the projection formula and the definition of degree over different
base fields. A complementary cycle on $X_{\bar k}$ is defined over some
finite extension $L/k$, so its pairing with $\alpha_{\bar k}$ is also zero.
The assumed algebraically closed case supplies an $N$ with
$\alpha_{\bar k}^{\boxtimes N}=0$.

A rational equivalence witnessing this vanishing uses finitely many
subvarieties, coefficients, and rational functions. All descend to a finite
extension $L/k$, after enlarging it if necessary. Thus
$\alpha_L^{\boxtimes N}=0$. For the finite flat map
$p_N:(X^N)_L\to X^N$,
\[
       (p_N)_*p_N^*(\alpha^{\boxtimes N})
                          =[L:k]\alpha^{\boxtimes N}=0.
\]
Division by $[L:k]$ in the rational Chow group proves the assertion.
Finite extensions need not be separable: finite flat pullback and proper
pushforward have the same degree identity. No characteristic-zero
hypothesis has entered. Components can be treated in each pure dimension,
as in the original statement. \hfill$\square$

This proves that passing to an algebraic closure is a legitimate reduction,
not an unannounced restriction. It does not prove the algebraically closed
case, nor does it give a way to decide whether the rational equivalence
over $\bar k$ exists.

\paragraph{Attempt: a computable exponent for a certified class.}
Suppose that over a finite extension $L/k$ one has an actual identity
\begin{equation}\label{eq182-certificate}
 \alpha_L=\sum_{i=1}^s(\Gamma_i)_*z_i,
 \qquad z_i\in\mathrm{CH}^1(C_i)_{\Q},\quad\deg z_i=0,
 \qquad \Gamma_i\in\mathrm{CH}^r(C_i\times X_L)_{\Q},
\end{equation}
where $C_i/L$ are smooth projective geometrically integral curves of
genus $g_i$. This is a \emph{certificate to be supplied}, not a consequence
of numerical triviality. It records the entire cycle identity, not merely
agreement of cohomology classes or intersection numbers.

For safety use $N_i=2g_i+1$. These integers are within the vanishing range
of Voevodsky's curve theorem \sref{182}{1}, Proposition 3.1, so
$z_i^{\boxtimes N_i}=0$. Functoriality of correspondences gives
\[
 ((\Gamma_i)_*z_i)^{\boxtimes N_i}
             =(\Gamma_i^{\boxtimes N_i})_*
                            (z_i^{\boxtimes N_i})=0.
\]
The following bookkeeping argument turns the separate vanishings into a
single explicit exponent:
\begin{equation}\label{eq182-exponent}
           \alpha^{\boxtimes N}=0,
             \qquad N=1+\sum_{i=1}^s(N_i-1)
                         =1+2\sum_{i=1}^s g_i.
\end{equation}
Indeed, expand the external power of the sum as a sum indexed by ordered
lists of $N$ indices. Some index $i$ occurs at least $N_i$ times. After a
permutation of factors, that summand contains the vanishing external power
of $(\Gamma_i)_*z_i$ as a factor, and hence is zero. All sums and
permutations are finite; no infinite tensor completion is used. The
descent lemma proves the same exponent over $k$. If the sum is empty,
$\alpha_L=0$ and exponent one suffices.

This also proves directly the familiar closure of smash-nilpotent cycles
under finite addition, with an explicit bound; compare
\sref{182}{1}, Lemma 2.6. For example, a certificate using curves of genera
one and two supplies exponent seven. The formula is not asserted to be
optimal, and it is not a computable exponent from $X$ and the numerical
hypothesis alone: the certificate \eqref{eq182-certificate} is additional
input.

\paragraph{Attempt: weaken the certificate instead of confusing equivalences.}
A more plausible sufficient input for a general numerically trivial
$\alpha$ is
\begin{equation}\label{eq182-intermediate}
 \text{for some }m\ge1,\quad
             \alpha^{\boxtimes m}\text{ is algebraically trivial on }X^m.
\end{equation}
Voevodsky's theorem applied on $X^m$ then gives an $N$ such that
\[
 (\alpha^{\boxtimes m})^{\boxtimes N}
                    =\alpha^{\boxtimes mN}=0.
\]
Conversely, smash-nilpotence gives \eqref{eq182-intermediate}, since a zero
Chow class is algebraically trivial. Thus, using the established theorem,
\eqref{eq182-intermediate} is an equivalent intermediate target. This
reduction is already contemplated in \sref{182}{1}, Section 4; it is
stated here to identify exactly what the proposed curve-certificate
strategy would have to produce.

Producing the certificate at $m=1$ would ask for more than this reduction.
The definitions of numerical and algebraic equivalence do not identify
them. Passing to a finite extension, or decomposing a motive into pieces
with nilpotent endomorphisms, has not constructed
\eqref{eq182-intermediate}. The next explicit test shows why one tempting
substitute for that construction is invalid.

\paragraph{Stress test: composition-square-zero is not tensor-nilpotent.}
Let $X=\PP^1_k\times\PP^1_k$. Denote the two divisor classes by $h_1,h_2$,
so $h_1^2=h_2^2=0$ and $\deg(h_1h_2)=1$. The degree-zero correspondence
\[
                 r=h_1\boxtimes h_1
                         \in\mathrm{CH}^2(X\times X)_{\Q}
\]
is nonzero. The composition formula for external-product correspondences
gives
\[
                    r\circ r=\deg(h_1^2)r=0.
\]
Nevertheless, for every $N\ge1$,
\begin{equation}\label{eq182-notensor}
 \deg\left(r^{\boxtimes N}\cdot
                      (h_2\boxtimes h_2)^{\boxtimes N}\right)=1.
\end{equation}
Thus none of the external powers of $r$ vanish. The equality follows by
factoring the intersection number over all $2N$ copies of $X$; each factor
is $\deg(h_1h_2)=1$. In particular $r$ is not numerically trivial, so it is
\emph{not} a counterexample to the target. It is a counterexample to the
auxiliary claim that composition nilpotence alone implies smash-nilpotence.
Any use of a Kimura-type nilpotence theorem would require an extra argument
exploiting its hypotheses, rather than just changing the multiplication
operation in its conclusion.

There is also no useful intersection-power shortcut: a positive-codimension
cycle has sufficiently high intersection powers zero on a fixed finite-
dimensional variety for dimensional reasons. Its external powers live on
varieties of growing dimension, so that automatic vanishing says nothing
about \eqref{eq182-intermediate}.

\paragraph{Stress test: dependence on cohomological standard conjectures.}
In any Weil cohomology theory with its K\"unneth identification, external
products satisfy
\[
 \operatorname{cl}(\alpha^{\boxtimes N})
                      =\operatorname{cl}(\alpha)^{\otimes N}.
\]
A nonzero vector over a field has every tensor power nonzero: choose a
linear functional not vanishing on it and take its tensor power. Hence
smash-nilpotence forces homological triviality. A universal proof of the
selected numerical-to-smash implication would therefore imply that
numerical and homological equivalence agree in these realizations, the
standard-conjecture consequence discussed in \sref{182}{1}, Section 4.
This is a warning about the depth of the missing input, not a reason to
abandon it or a conditional assumption quietly used in the descent proof.

\paragraph{Outcome.}
\textbf{Outcome: Conditional reduction.}

The attempt proves descent without separability restrictions, an explicit
external-power bound for supplied curve-correspondence certificates, and
the implication from an algebraically trivial external power to the target.
It also gives an exact correspondence showing the failure of a
composition-nilpotence shortcut. The underlying nilpotence theorem and
intermediate reduction are established literature; no independent novelty
claim is made for them or for the bookkeeping above.

The missing result is \eqref{eq182-intermediate} for an arbitrary
numerically trivial rational cycle, or another construction of an external
rational-equivalence certificate. Neither the descent nor the tropical
analogue supplies it. All-field scope is retained, but the full conjecture
has not been resolved.
% END RESEARCH ADDITION ATTEMPT-182
\subsection{Sources}
\begin{itemize}\raggedright
\item[\textnormal{[S1]}]\hypertarget{source:182:1}{} Vladimir Voevodsky, A Nilpotence Theorem for Cycles Algebraically Equivalent to Zero, IMRN1995 no.4,187-198; published IAS snapshot,13pages.
\url{https://www.math.ias.edu/vladimir/sites/math.ias.edu.vladimir/files/Nilpotence_theorem_published.pdf}.
{\small\textit{Catalogue formulation source.}}
% BEGIN RESEARCH ADDITION SOURCES-182
\item[\textnormal{[E1]}]\hypertarget{extra:182:1}{} Alexia Corradini,
\emph{Algebraically trivial tropical cycles are smash-nilpotent},
arXiv:2608.21988v1, 22 August 2026, abstract and stated tropical scope.
\url{https://arxiv.org/abs/2608.21988v1}.
{\small\textit{Consulted 27 September 2026. Adjacent-category result;
no algebraic Chow-group theorem is inferred from it.}}
% END RESEARCH ADDITION SOURCES-182
\end{itemize}

\end{document}
